\documentclass[12pt, letterpaper]{article}
\date{\today}
\usepackage[margin=1in]{geometry}
\usepackage{amsmath}
\usepackage{hyperref}
\usepackage{cancel}
\usepackage{amssymb}
\usepackage{fancyhdr}
\usepackage{pgfplots}
\usepackage{booktabs}
\usepackage{pifont}
\usepackage{amsthm,latexsym,amsfonts,graphicx,epsfig,comment}
\pgfplotsset{compat=1.16}
\usepackage{xcolor}
\usepackage{tikz}
\usetikzlibrary{shapes.geometric}
\usetikzlibrary{arrows.meta,arrows}
\newcommand{\Z}{\mathbb{Z}}
\newcommand{\N}{\mathbb{N}}
\newcommand{\R}{\mathbb{R}}
\newcommand{\Q}{\mathbb{Q}}
\newcommand{\C}{\mathbb{C}}
\newcommand{\F}{\mathbb{F}}

\newcommand{\Po}{\mathcal{P}}
\newcommand{\Pro}{\mathbb{P}}
\author{Alex Valentino}
\title{501 homework}
\pagestyle{fancy}
\renewcommand{\headrulewidth}{0pt}
\renewcommand{\footrulewidth}{0pt}
\fancyhf{}
\rhead{
	Homework 5\\
	501
}
\lhead{
	Alex Valentino\\
}
\begin{document}
\begin{enumerate}
	\item[6] Let $\mu$ be the Lesbesgue measure on $\R$ and let $E \subset \R$ be a Lesbesgue measurable set with 
	$\mu(E) > 0$.  We want to show that the set $E - E = \{ x - y: x,y \in E\}$
	contains an open interval centered at the origin.  Since $\mu(E) > 0$ then we can
	write $\mu(E) = \sum_{n = -\infty}^\infty \mu(E \cap (n,n+1))$ and be assured that 
	there exists $n'$ such that $\mu(E \cap [n',n'+1]) > 0$.  If we consider $E - n$ 
	then we can apply problem 5 and find an open interval $I \subseteq (n,n+1)$ such that $\mu(E \cap I) > r \mu(I)$ for all $r \in (0,1)$.  Note that this implies that 
	$\mu(E \cap I) = \mu(I)$.  Therefore we know for arbitrary $\epsilon > 0$ 
	that $\mu(E \cap [n,n+1] \Delta I) < \epsilon$.  We claim that the function 
	$\phi(x) = \mu((E\cap I)\cap (E \cap I + x))$ is continuous.  Note that for the 
	function $\psi(x) = \mu(E \cap I \cap I + x)$ we can compute it exactly, given that 
	there exists $n < a < b < n+1$ such that $I = (a,b)$ then we have that 
	$\psi(x) - \psi(y) = \mu( E \cap I \cap I + x) - \mu(E \cap I \cap I + y)  b - a - |x| + b - a + |y| = |y| - |x|$.  where $x,y \in (a-b, b -a)$.  Furthermore, if we consider 
	$\phi(x) - \psi(x) = \mu((E\cap I)\cap (E \cap I + x)) - \mu(I \cap I + x)$, 
	since we have that 
	\begin{align*}
	\mu((E\cap I)\cap (E \cap I + x) \Delta (I \cap I + x)) &= \mu( (I \cap I + x) \backslash (E\cap I)\cap (E \cap I + x))\\
	&= \mu(
	\end{align*}
	I didn't have complete time, but if I could sus out the convergence of the 
	$\phi$ then from that one can take an $\epsilon$ neighborhood around some $\psi$ 
	which is non-zero, then from there you can look at the set it gives out in $\psi^{-1}$.  Since this will be an open set you can then subtract every element by 
	every element within the delta neighborhood, which has the effect of translating
	the neighborhood to have an edge on the origin, include the origin, then reflect
	it accross the axis to reflect sign flips of the elements subtracting.
	\item[7] Let $\mathcal{F}$ be a countable family of lesbesgue measurable functions, $E$ be a lesbesgue measurable set with $\mu(E) < \infty$ with the property for each 
	$x \in E$ there exists $M_x < \infty$ such that for all $f \in \mathcal{F}, |f(x)| \leq M_x$, and $\epsilon >0$.  We want to construct a closed set $F \subset E$ such that there exists $M > 0$ that for all $x \in F, f \in \mathcal{F}, f(x) \leq M$ with 
	$\mu(E \cap F^c) \leq \epsilon$.  Let 
	$f' = \sup_{f \in \mathcal{F}} f$.  Note that $f'$ is a Lesbesgue measurable 
	function.  Note that since every point $x \in E$ every function in $\mathcal{F}$
	has a finite pointwise bound.  Therefore $f'$ takes on $\R$ values on $E$.  
	Therefore $f'^{-1}((-\infty,\infty)) \cap E = E$, as every point in $E$ is 
	eventually included.  Therefore there exists $\infty > M > 0$ such that 
	$\mu(f'^{-1}([-M,M]) \cap E) \geq \mu(E) - \epsilon/2$.  Furthermore by the inner 
	regularity of $\mu$ we have that there exists $F \subset f'^{-1}([-M,M]) \cap E$
	such that $\mu(F) \geq \mu(f'^{-1}([-M,M]) \cap E) - \epsilon/2$.  Therefore we have that
	\begin{align*}
		\mu(E) &= \mu(F) + \mu(E \cap F^c)\\
		&\geq \mu(E) - \epsilon + \mu(E \cap F^c)\\
		\epsilon &\geq \mu(E \cap F^c)
	\end{align*}
	Additionally, by taking the preimage of $[-M,M]$ of the sup of the functions of 
	$\mathcal{F}$ the uniform boundedness is implied.  
	\item[8] Let $E$ be a non-lesbesgue measurable set and let $A$ be a lesbesgue measurable set of measure 0.  We claim that $E \cap A^c$ is not Lesbesgue measurable.
	Suppose not.  Then $E \cap A^c$ is lesbesgue measurble  Additionally, since $E$ is
	not Lesbesgue measurable then there exists $C \subset \R$ such that 
	$\mu^*(C) < \mu^*(E \cap C) + \mu^*(E^c \cap C)$, where $\mu^*$ is the outer lesbesgue measure.  Furthermore, we note that 
	$E \cap A$ must be Lesbesgue measurable since $E \cap A \subseteq A$, and therefore
	$0 \leq \mu^*(E \cap A) \leq \mu(A) = 0$, which by the definition of the completeness
	of measure spaces implies that $E\cap A$ is Lebesgue measurable (a similar 
	argument works for $\mu^*(E^c \cap A)$.  Since $A \cap E, A \cap E^c, A^c \cap E$
	are all measurable then their union is measurable, $A \cup E$, and the 
	complement of their union is measurable, $A^c \cap E^c$.  
	
	Therefore 
	we have that 
	\begin{align*}
	\mu^*(C) &< \mu^*(C \cap E) + \mu^*(C \cap E^c)\\
	&= \mu^*(C \cap E \cap A) +\mu^*(C \cap E \cap A^c) + \mu^*(C \cap E^c \cap A)
	+\mu^*(C \cap E^c \cap A^c)\\
	&= \mu^*(C \cap (E \cup A)) + \mu^*(C \cap E^c \cap A^c)\\
	&= \mu^*c(C).
	\end{align*}
	Note that this is a contradiction.  
	\item[9] Let $A_0 \subset A_1 \subset \cdots \subset A_n \subset \cdots [0,1]$.
	We want to show that $\lim \mu^*(A_n) = \mu^*(\bigcup_{n=1}^\infty A_n)$.
	Note by the outer regularity of $\mu^*$ there exists $B_n$ open such that 
	$A_n \subset B_n$ and $\mu^*(B_n) \leq \mu^*(A_n) + 2^{-n}$, and 
	$C_n = \cap_{m=n}^\infty B_m$.  Since each $C_n$ is a measurable set then we 
	have that $\lim \mu(C_n) = \mu(\cup_{n=1}^\infty C_n)$.
	Note that $\mu(\cup_{n=1}^\infty C_n) \leq \sum_{n=1}^\infty \mu(C_n)$.
	Furthermore, if we consider that (I ran out of time, sorry).
	\item[10] Let $\mu$ be the Lebesgue measure on $\R$, and let $0 \leq a < b \leq 1$.
	We will construct an open set $U \subset [0,1]$ such that $\mu(A) = a, \mu(B) = b$.
	If we restrict ourselves to the interval $[0,b]$, and choose an $\alpha \geq 1$ 
	such that $\frac{a}{b} < x', 1 - 2x'^\alpha - x', x' >0$.  Then we begin to 
	construct the set $C$ as follows: for $C_1$, we remove the middle $r$ portion of 
	$[0,b]$ and we're left with $[0, \frac{r}{2}] \cup [\frac{r}{2} - b, b]$.  Then for 
	$C_2$ we remove the middle $r^(\alpha + 1)$ terms from each remaining interval, 
	divinding them in two.  At the $n$th step we remove $2^{n-1}r^{1+\alpha(n-1)}$ from 
	the collective remaining intervals.  The final set we end up with is
	$C = \cap_{j=1}^\infty C_j$.  Note that $C$ is closed.  For our desired set, we take
	$U = C^c$.  Note that $\mu(U) = b - b\sum_{i=1}^\infty 2^{n-1} r^{1 + \alpha(n-1)}
	= b\frac{1 - 2r^\alpha - r}{1 - 2r^\alpha}$.  Note that since $\frac{a}{b} < 1 - 2x'^\alpha - x'$ then we have that the numerator is not 0, and ensure that the function
	will be positive.  Then we can choose a positive $r$ such that $a = \frac{1 - 2r^\alpha - r}{1 - 2r^\alpha}$, since the area function is continuous on $[0,x']$, and
	by the the intermediate value theorem since $r = 0$ corresponds to 1 and $r=x'$ 
	corresponds to 0, and in the region our function is positive.  Therefore 
	$\mu(U) = b \frac{a}{b} = a$.  Furthermore, $\mu(\bar{U}) = b$ since $U$ is the 
	interval excluding well chosen limit points.  
\end{enumerate}
\end{document}
